\section{Limitations and implications}
The results describe population identification on a finite ordered rating scale.
They assume a specified conditional target distribution for each item.
In practice, a small number of ratings estimates that distribution imprecisely.
Changes in the rater population or the meaning of rating categories change the target itself.
One observed rating supports an unbiased individual-rating CE or RPS observation under the assumed sampling law, but it does not reveal the full conditional distribution for that item.

The geometry of unrestricted distributions does not automatically describe a restricted neural model.
Parameter sharing, regularization, initialization, and optimization may choose a particular shape even when the objective does not identify it.
Conversely, strict properness does not remove misspecification or optimization error.
The linear-model result needs covariance and feasibility conditions, and the exponential-family result needs its stated sufficient statistic and positive variance.
Global identification requires checking the full feasible affine-copy set; local derivative conditions provide only the stated local guarantee.

The exact illustration is a mathematical construction with no trained judge or sampled human ratings.
It establishes possibility and reproduces a loss difference exactly.
It does not estimate the prevalence of such differences or compare real-world training methods.
The asymptotic CE calculation concerns regular, interior, well-specified iid maximum likelihood with additional expectation control.
It does not predict a universal sample size at which CE-trained judges surpass correlation-trained judges.

Our main results concern Pearson correlation of conditional means.
Spearman and Kendall coefficients computed only from those means also discard equal-mean shape differences, but their equality conditions and gradients differ.
Probability-profile correlation is a separate objective with the ambiguity in Proposition~\ref{prop:profiles}.
For applications that require only rank association, distribution recovery can exceed the task's requirements.
For applications that interpret a judge's output as a rating distribution, individual-rating proper scores directly target that distribution.
Correlation remains a useful association diagnostic when its target and undefined constant-vector cases are reported explicitly.
